\documentclass[11pt,a4paper]{article}

\usepackage[utf8]{inputenc}
\usepackage[T1]{fontenc}
\usepackage{amsmath, amssymb, amsthm}
\usepackage{geometry}
\usepackage{hyperref}
\usepackage{booktabs}
\usepackage{array}
\usepackage{multirow}
\usepackage{graphicx}
\usepackage{caption}
\usepackage{float}
\usepackage{enumitem}
\usepackage{mathtools}

\geometry{a4paper, margin=1in}

\theoremstyle{plain}
\newtheorem{theorem}{Theorem}
\newtheorem{proposition}[theorem]{Proposition}
\newtheorem{lemma}[theorem]{Lemma}
\newtheorem{corollary}[theorem]{Corollary}
\newtheorem{conjecture}[theorem]{Conjecture}
\theoremstyle{definition}
\newtheorem{definition}[theorem]{Definition}
\theoremstyle{remark}
\newtheorem{remark}[theorem]{Remark}

\DeclareMathOperator{\li}{li}

\title{Prime values and consecutive-prime clusters of \texorpdfstring{$n^5-(n-1)^5$}{n\^{}5-(n-1)\^{}5}: \\[4pt] Computation, heuristics, and large probable prime generation}

\author{Ruqing Chen\thanks{GUT Geoservice Inc., Montr\'eal, QC, Canada.  Email: \href{mailto:ruqing@hotmail.com}{\texttt{ruqing@hotmail.com}}}}

\date{}

\begin{document}

\maketitle

\begin{abstract}
We report on an exhaustive computational search for prime values of the quintic difference polynomial $Q(n) = n^5 - (n-1)^5$ over the range $2 \le n \le 10^8$, yielding $5{,}179{,}467$ primes.
We derive the Bateman--Horn constant $C_Q = 3.678113109 \pm 3 \times 10^{-9}$ from the factorization structure of $Q$ over the cyclotomic field $\mathbb{Q}(\zeta_5)$, and verify that the observed prime counts agree with the prediction $\pi_Q(N) \sim C_Q \int_2^N dt/\!\log Q(t)$ to within $0.01\%$ at $N = 5 \times 10^7$.
We analyze the distribution of consecutive-prime $k$-tuplets (sequences $n, n{+}1, \ldots, n{+}k{-}1$ for which $Q(n), Q(n{+}1), \ldots, Q(n{+}k{-}1)$ are all prime) and prove that $7$-tuplets cannot exist, while $6$-tuplets can and do occur.
Extended searches up to $n \approx 1.9 \times 10^9$ have located $25$ sextuplets, $458$ quintuplets, and $11{,}063$ quadruplets.
Hardy--Littlewood-type singular series for these $k$-tuplets are computed and shown to match observed counts.
Finally, we describe a three-stage pipeline (C++ modular sieve, target construction, PFGW probable-prime testing) that has produced $175$ probable primes $Q(n)$ with $\sim\! 10{,}000$ digits, two with $\sim\! 50{,}000$ digits, one with $\sim\! 100{,}000$ digits, and one with $\sim\! 200{,}000$ digits. One of the $10{,}000$-digit candidates has been rigorously certified prime via ECPP, yielding a $9{,}997$-digit proven prime of the form $n^5-(n-1)^5$.
\end{abstract}

\bigskip
\noindent \textbf{2020 Mathematics Subject Classification.} 11A41, 11N32, 11Y11, 11Y35.

\noindent \textbf{Key words.} Prime-generating polynomials, Bateman--Horn conjecture, prime tuplets, cyclotomic polynomials, probable primes, large prime search.

% ====================================================================
\section{Introduction}
\label{sec:intro}
% ====================================================================

For a fixed polynomial $f \in \mathbb{Z}[x]$, the question of how often $f(n)$ takes prime values as $n$ ranges over the positive integers is classical and largely open. Aside from the linear case $f(n) = an + b$ (Dirichlet's theorem on primes in arithmetic progressions), no polynomial of degree $\ge 2$ has been proved to represent infinitely many primes. Nevertheless, the Bateman--Horn conjecture~\cite{bateman1962} provides a precise quantitative heuristic, and extensive numerical evidence supports it for many families; see Sorenson~\cite{sorenson2024} for a recent survey.

In this paper we study the polynomial
\begin{equation}\label{eq:Qdef}
Q(n) \;=\; n^5 - (n-1)^5 \;=\; 5n^4 - 10n^3 + 10n^2 - 5n + 1,
\end{equation}
which arises naturally as the fifth-power analogue of the well-known case $n^2 - (n-1)^2 = 2n - 1$, whose prime values are simply the odd primes. We note the symmetry $Q(n) = Q(1-n)$, which follows from the substitution $n \mapsto 1-n$ in \eqref{eq:Qdef}; consequently, the study of prime values for negative $n$ reduces trivially to the positive case. The cubic case $n^3-(n-1)^3 = 3n^2-3n+1$ has been investigated as a source of large primes by Dubner and others~\cite{dubner2002}, but the quintic case appears not to have received systematic attention.

Two features make $Q(n)$ particularly attractive:
\begin{enumerate}[label=(\roman*)]
\item \textbf{Arithmetic rigidity.} The factorization of $Q$ is governed by the fifth cyclotomic polynomial $\Phi_5$, which forces every prime divisor of $Q(n)$ to be congruent to $1 \pmod{5}$ (see Corollary~\ref{cor:divisibility}). This eliminates three-quarters of all primes from the sieve, producing a large Bateman--Horn constant $C_Q \approx 3.678$ --- meaning $Q(n)$ is roughly $3.68$ times as likely to be prime as a ``random'' integer of the same size.

\item \textbf{Tuplet clustering.} The same arithmetic rigidity that boosts overall prime density also permits long runs of consecutive $n$-values for which $Q(n)$ is simultaneously prime. We call a maximal run of $k$ consecutive integers $\{n, n{+}1, \ldots, n{+}k{-}1\}$ all yielding primes a \emph{$k$-tuplet}. We prove that $k \le 6$ is a hard upper bound imposed by the roots of $Q$ modulo~$11$, and exhibit data showing that this bound is achieved.
\end{enumerate}

The paper is organized as follows. Section~\ref{sec:arithmetic} develops the arithmetic of $Q(n)$ modulo primes, establishing the root structure and admissibility constraints. Section~\ref{sec:BH} derives the Bateman--Horn prediction and compares it with our computational data. Section~\ref{sec:gaps} analyzes the prime gap distribution. Section~\ref{sec:tuplets} treats $k$-tuplets and their Hardy--Littlewood heuristics. Section~\ref{sec:large} describes the algorithmic pipeline for generating large probable primes of the form $Q(n)$, and Section~\ref{sec:conclusion} concludes with open problems.


% ====================================================================
\section{Arithmetic of $Q(n)$ modulo primes}
\label{sec:arithmetic}
% ====================================================================

\subsection{Cyclotomic structure}

The identity
\begin{equation}\label{eq:factor}
n^5 - (n-1)^5 \;=\; (n - (n-1)) \cdot \sum_{j=0}^{4} n^j (n-1)^{4-j} \;=\; \sum_{j=0}^{4} n^j (n-1)^{4-j}
\end{equation}
shows that $Q(n)$ is the norm form arising from the factorization of $x^5 - 1$. More precisely, setting $x = n/(n-1)$, we have
\begin{equation}\label{eq:Phi5}
Q(n) \;=\; (n-1)^4 \,\Phi_5\!\left(\frac{n}{n-1}\right),
\end{equation}
where $\Phi_5(x) = x^4 + x^3 + x^2 + x + 1$ is the fifth cyclotomic polynomial. The splitting behavior of $\Phi_5$ modulo a prime $p$ is determined by the order of $p$ in $(\mathbb{Z}/5\mathbb{Z})^*$:

\begin{proposition}\label{prop:roots}
Let $p$ be a prime and let $\omega_Q(p)$ denote the number of solutions $n \in \{0, 1, \ldots, p-1\}$ to $Q(n) \equiv 0 \pmod{p}$. Then:
\begin{enumerate}[label=\emph{(\alph*)}]
\item If $p = 5$, then $\omega_Q(5) = 0$. In fact, $Q(n) \equiv 1 \pmod{5}$ for all $n$.
\item If $p \equiv 1 \pmod{5}$, then $\omega_Q(p) = 4$.
\item If $p \not\equiv 0, 1 \pmod{5}$, then $\omega_Q(p) = 0$.
\end{enumerate}
\end{proposition}

\begin{proof}
(a) By Fermat's little theorem, $n^5 \equiv n \pmod{5}$ for all $n$, so $Q(n) = n^5 - (n-1)^5 \equiv n - (n-1) = 1 \pmod{5}$.

(b)--(c) By \eqref{eq:Phi5}, for $p \nmid n-1$ we have $Q(n) \equiv 0 \pmod{p}$ if and only if $\Phi_5(n/(n-1)) \equiv 0 \pmod{p}$, i.e., $n/(n-1)$ is a primitive fifth root of unity modulo~$p$. When $p \equiv 1 \pmod{5}$, the group $(\mathbb{Z}/p\mathbb{Z})^*$ contains exactly four primitive fifth roots of unity; each yields a unique $n \equiv \zeta/(\zeta - 1) \pmod{p}$ with $\zeta \ne 1$. When $p \not\equiv 1 \pmod{5}$ and $p \ne 5$, there are no primitive fifth roots of unity in $(\mathbb{Z}/p\mathbb{Z})^*$, so $\Phi_5$ has no roots. One verifies separately that $n \equiv 1 \pmod{p}$ does not yield $Q(n) \equiv 0 \pmod{p}$ (since $Q(1) = 1$).
\end{proof}

\begin{corollary}\label{cor:divisibility}
Every prime factor of $Q(n)$ is congruent to $1 \pmod{5}$. In particular, $Q(n)$ is never divisible by $2, 3, 7,$ or any prime $p \not\equiv 1 \pmod{5}$.
\end{corollary}

\begin{proof}
If $p \mid (n-1)$ then $Q(n) \equiv 1 \pmod{p}$, so $p \nmid Q(n)$. For $p \nmid (n-1)$, Proposition~\ref{prop:roots} gives the result.
\end{proof}

The smallest prime $\equiv 1 \pmod{5}$ is $p = 11$. Its four roots are:
\begin{equation}\label{eq:roots11}
Q(n) \equiv 0 \pmod{11} \quad\Longleftrightarrow\quad n \equiv 4, 5, 7, 8 \pmod{11}.
\end{equation}
Thus $Q(n)$ can be prime only when $n \bmod 11 \in \{0, 1, 2, 3, 6, 9, 10\}$, a set of $7$ residues.

Table~\ref{tab:roots} lists the roots for the first several primes $p \equiv 1 \pmod{5}$.

\begin{table}[ht]
\centering
\caption{Roots of $Q(n) \equiv 0 \pmod{p}$ for primes $p \equiv 1 \pmod{5}$}
\label{tab:roots}
\smallskip
\begin{tabular}{@{}rllc@{}}
\toprule
$p$ & Roots $\{r : Q(r) \equiv 0\}$ & Largest consecutive safe run & $\#$ safe residues \\
\midrule
$11$  & $\{4, 5, 7, 8\}$     & $6$ \;\;$(9,10,0,1,2,3)$ & $7$ \\
$31$  & $\{2, 10, 22, 30\}$   & $11$  & $27$ \\
$41$  & $\{9, 12, 30, 33\}$   & $17$  & $37$ \\
$61$  & $\{16, 24, 38, 46\}$  & $30$  & $57$ \\
$71$  & $\{4, 19, 53, 68\}$   & $33$  & $67$ \\
$101$ & $\{27, 29, 73, 75\}$  & $52$  & $97$ \\
\bottomrule
\end{tabular}\\[4pt]
\small Note: For every prime $p \equiv 1 \pmod{5}$, the number of safe residues is exactly $p - 4$.
\end{table}

\subsection{Admissibility and the maximal tuplet length}

\begin{theorem}\label{thm:7tuplet}
There is no integer $n$ such that $Q(n), Q(n{+}1), \ldots, Q(n{+}6)$ are all prime. That is, $7$-tuplets for $Q$ do not exist. Moreover, $6$ is the maximal tuplet length, and this bound is achieved.
\end{theorem}

\begin{proof}
By \eqref{eq:roots11}, the set of ``forbidden'' residues modulo~$11$ is $F = \{4, 5, 7, 8\}$. The complementary ``safe'' residues are $S = \{0, 1, 2, 3, 6, 9, 10\}$. The maximal length of a run of consecutive integers drawn entirely from $S$ is determined by the largest gap in $F$ when the residues $\{0, \ldots, 10\}$ are arranged cyclically. The gaps between consecutive elements of $F$ are $5-4=1$, $7-5=2$, $8-7=1$, and $(4+11)-8 = 7$. The largest gap is $7$, corresponding to the run $9, 10, 0, 1, 2, 3$ of length $6$. Hence any $7$ consecutive integers must contain at least one element of $F$, making $Q$ divisible by $11$ at that point.

Conversely, computations described in Section~\ref{sec:tuplets} exhibit $25$ sextuplets in the range $n \le 1.9 \times 10^9$, the smallest at $n = 10{,}291{,}136$.
\end{proof}


% ====================================================================
\section{The Bateman--Horn prediction}
\label{sec:BH}
% ====================================================================

\subsection{Statement and constant computation}

The Bateman--Horn conjecture~\cite{bateman1962} predicts that for an irreducible polynomial $f \in \mathbb{Z}[x]$ with positive leading coefficient and no fixed prime divisor,
\begin{equation}\label{eq:BH}
\pi_f(N) \;:=\; \#\{n \le N : f(n) \text{ is prime}\} \;\sim\; C_f \int_2^N \frac{dt}{\log f(t)},
\end{equation}
where the product
\begin{equation}\label{eq:BHconst}
C_f \;=\; \prod_{p\;\text{prime}} \frac{1 - \omega_f(p)/p}{1 - 1/p}
\end{equation}
converges to a positive constant (here $\omega_f(p)$ denotes the number of roots of $f$ modulo~$p$).

For $f = Q$, Proposition~\ref{prop:roots} gives:
\begin{equation}\label{eq:CQ}
C_Q \;=\; \frac{5}{4}\;\cdot\; \prod_{\substack{p \neq 5 \\ p \not\equiv 1\,(5)}} \frac{p}{p-1} \;\cdot\; \prod_{p \equiv 1\,(5)} \frac{p-4}{p-1}\,.
\end{equation}
The convergence follows from Dirichlet's theorem on primes in arithmetic progressions: primes $p \equiv 1 \pmod{5}$ have natural density $1/4$ among all primes, so the average of $\omega_Q(p)$ is $4 \cdot \frac{1}{4} + 0 \cdot \frac{3}{4} = 1$, and the logarithmic contributions
\[
\sum_p \frac{\omega_Q(p) - 1}{p}
\]
converge by partial summation and the Siegel--Walfisz theorem (see \cite[Ch.~17]{iwaniec2004} for the general theory).

We evaluate \eqref{eq:CQ} numerically by computing the partial Euler product $C_Q(P) = \prod_{p \le P} (1 - \omega_Q(p)/p)/(1 - 1/p)$ over all primes $p \le P$.  The tail error admits a precise bound: writing $f(p) = \log[(1 - \omega_Q(p)/p)/(1 - 1/p)]$ and using $\log(1-x) = -x - x^2/2 + O(x^3)$, we obtain
\begin{equation}\label{eq:tailbound}
f(p) = \frac{1 - \omega_Q(p)}{p} + \frac{1 - \bigl(\omega_Q(p)\bigr)^2}{2p^2} + O(p^{-3}).
\end{equation}
The first-order terms $\sum_{p > P}(1 - \omega_Q(p))/p$ cancel by Dirichlet's theorem, since $\omega_Q$ averages to~$1$ over primes in each residue class modulo~$5$.  The dominant tail contribution comes from the second-order terms: for $p \equiv 1 \pmod{5}$ (where $\omega_Q = 4$), the coefficient is $(1 - 16)/(2p^2) = -\frac{15}{2p^2}$, while for $p \not\equiv 1 \pmod{5}$ (where $\omega_Q = 0$), it is $1/(2p^2)$.  By the prime-number theorem for arithmetic progressions (each class carries density~$1/4$),
\[
\log C_Q - \log C_Q(P) = \Big({-}\tfrac{15}{8} + \tfrac{3}{8}\Big)\sum_{p > P}\frac{1}{p^2} + O\Big(\sum_{p>P}\frac{1}{p^3}\Big) = -\frac{3}{2}\sum_{p > P}\frac{1}{p^2} + O(P^{-2}\log^{-1}\!P).
\]
By partial summation with $\pi(t) \sim t/\!\log t$, one has $\sum_{p > P}1/p^2 \sim 1/(P\log P)$, so
\[
|\log C_Q - \log C_Q(P)| \;\sim\; \frac{3}{2P\log P}\,.
\]
Extending the computation to all primes $p \le 10^8$ yields (Table~\ref{tab:Cconvergence})
\begin{equation}\label{eq:Cvalue}
C_Q \;=\; 3.678113109 \pm 3 \times 10^{-9},
\end{equation}
where the error bound follows from the tail estimate above with $P = 10^8$ and $\log P \approx 18.4$.  Higher precision could be obtained by expressing $C_Q$ in closed form via Dirichlet $L$-functions $L(1, \chi)$ for the nontrivial characters $\chi \bmod 5$ (cf.~\cite{cohen2007,languasco2010}).  The relevant $L$-values are:
\begin{align*}
L(1,\chi_2) &= \frac{2\log\varphi}{\sqrt{5}} = 0.43040894096\ldots \quad\text{(class number formula, $\varphi = \tfrac{1+\sqrt{5}}{2}$)},\\
|L(1,\chi_1)|^2 &= \frac{[\psi(\tfrac{1}{5})-\psi(\tfrac{4}{5})]^2 + [\psi(\tfrac{2}{5})-\psi(\tfrac{3}{5})]^2}{25} = 0.78956835209\ldots,
\end{align*}
where $\psi$ is the digamma function and $\chi_1$ is the primitive character of order~$4$ modulo~$5$.  Full evaluation of $C_Q$ to arbitrary precision via these values (combined with an accelerated Euler product over small primes) is straightforward in a computer algebra system such as PARI/GP, and would yield a value consistent with \eqref{eq:Cvalue}.

\begin{table}[ht]
\centering
\caption{Convergence of the partial Euler product for $C_Q$}
\label{tab:Cconvergence}
\smallskip
\begin{tabular}{@{}rc@{}}
\toprule
Primes $p \le$ & Partial product \\
\midrule
$10^2$ & $3.5557$ \\
$10^3$ & $3.6746$ \\
$10^4$ & $3.6715$ \\
$10^5$ & $3.6777$ \\
$10^6$ & $3.6777$ \\
$10^7$ & $3.6781$ \\
$10^8$ & $3.678113109$ \\
\bottomrule
\end{tabular}
\end{table}


\subsection{Comparison with data}

We computed all $n \le 10^8$ for which $Q(n)$ is prime by a segmented sieve exploiting the root structure of Proposition~\ref{prop:roots}: for each prime $p \equiv 1 \pmod{5}$ up to a suitable bound, we mark the four residue classes where $Q(n) \equiv 0 \pmod{p}$, then test survivors for primality (see~\cite{galway2004} for related sieve methodology). The search found $\pi_Q(10^8) = 5{,}179{,}467$ primes. Table~\ref{tab:BH} compares the observed counts with the Bateman--Horn prediction $C_Q \int_2^N dt/\!\log Q(t)$ using $C_Q = 3.6781$.

\begin{table}[ht]
\centering
\caption{Comparison of $\pi_Q(N)$ with Bateman--Horn prediction}
\label{tab:BH}
\smallskip
\begin{tabular}{@{}rrrr@{}}
\toprule
$N$ & $\pi_Q(N)$ & $C_Q \int_2^N \frac{dt}{\log Q(t)}$ & Ratio \\
\midrule
$10^3$           &         $149$ &       $151.1$ & $0.9864$ \\
$10^4$           &       $1{,}101$ &     $1{,}088.6$ & $1.0114$ \\
$10^5$           &       $8{,}489$ &     $8{,}518.3$ & $0.9966$ \\
$10^6$           &      $69{,}977$ &    $70{,}061.2$ & $0.9988$ \\
$2\times 10^6$   &     $133{,}203$ &   $133{,}036.3$ & $1.0013$ \\
$5\times 10^6$   &     $312{,}522$ &   $311{,}772.9$ & $1.0024$ \\
$10^7$           &     $596{,}022$ &   $595{,}382.1$ & $1.0011$ \\
$2\times 10^7$   &   $1{,}139{,}867$ & $1{,}139{,}339.1$ & $1.0005$ \\
$5\times 10^7$   &   $2{,}694{,}860$ & $2{,}694{,}614.7$ & $1.0001$ \\
\bottomrule
\end{tabular}
\end{table}

The agreement is striking: at $N = 5 \times 10^7$, the observed count exceeds the prediction by only $245$ primes out of $2.69$ million, a relative error of $0.009\%$.


\subsection{Residue class uniformity}

As a further consistency check, Corollary~\ref{cor:divisibility} predicts that prime-generating values of $n$ should be uniformly distributed among the $7$ safe residue classes modulo~$11$. Table~\ref{tab:residues} confirms this prediction to high precision.

\begin{table}[ht]
\centering
\caption{Distribution of $n \pmod{11}$ among all $5{,}179{,}467$ primes $Q(n)$}
\label{tab:residues}
\smallskip
\begin{tabular}{@{}rrc@{}}
\toprule
$n \bmod 11$ & Count & Status \\
\midrule
$0$  & $738{,}642$ & safe \\
$1$  & $739{,}726$ & safe \\
$2$  & $740{,}210$ & safe \\
$3$  & $739{,}507$ & safe \\
$4$  & $0$ & root \\
$5$  & $0$ & root \\
$6$  & $740{,}038$ & safe \\
$7$  & $0$ & root \\
$8$  & $0$ & root \\
$9$  & $740{,}500$ & safe \\
$10$ & $740{,}844$ & safe \\
\bottomrule
\end{tabular}\\[4pt]
\small Expected per safe class: $739{,}924$. Maximum deviation: $0.12\%$.
\end{table}


% ====================================================================
\section{Prime gap distribution}
\label{sec:gaps}
% ====================================================================

Let $n_1 < n_2 < \cdots$ be the sequence of positive integers for which $Q(n_i)$ is prime, and set $g_i = n_{i+1} - n_i$. Over $5{,}179{,}466$ gaps in our dataset, we find a mean gap of $\bar{g} = 19.31$ with a maximum of $284$. (For comparison with gap statistics for ordinary primes, see Oliveira e~Silva et~al.~\cite{oliveira2014}.)

\subsection{Raw gap frequencies}

The gap distribution displays a pronounced mod-$11$ structure inherited from the forbidden residues \eqref{eq:roots11}. Table~\ref{tab:gapfreq} shows frequencies for small gaps.

\begin{table}[ht]
\centering
\caption{Gap frequencies $g_i = n_{i+1} - n_i$ for small values}
\label{tab:gapfreq}
\smallskip
\begin{tabular}{@{}rrr|rrr@{}}
\toprule
Gap & Count & Frequency & Gap & Count & Frequency \\
\midrule
$1$  & $288{,}450$ & $5.57\%$ & $7$  & $182{,}678$ & $3.53\%$ \\
$2$  & $218{,}017$ & $4.21\%$ & $8$  & $237{,}959$ & $4.59\%$ \\
$3$  & $290{,}431$ & $5.61\%$ & $9$  & $165{,}254$ & $3.19\%$ \\
$4$  & $196{,}226$ & $3.79\%$ & $10$ & $190{,}263$ & $3.67\%$ \\
$5$  & $139{,}104$ & $2.69\%$ & $11$ & $266{,}046$ & $5.14\%$ \\
$6$  & $141{,}793$ & $2.74\%$ & $22$ & $144{,}659$ & $2.79\%$ \\
\bottomrule
\end{tabular}
\end{table}

Several features are noteworthy. Gaps $g = 1$ and $g = 3$ are the most common, reflecting the pair structure of safe residues modulo~$11$: the safe set $\{0, 1, 2, 3, 6, 9, 10\}$ contains the consecutive blocks $\{0,1,2,3\}$ and $\{9,10\}$, favoring short gaps. The prominence of $g = 11$ (and multiples thereof) arises because stepping by $11$ preserves the residue class modulo~$11$.

\subsection{Normalized gaps and statistical testing}

To remove the slow variation in gap size, we normalize: $\hat{g}_i = g_i / E_i$ where $E_i = (4\log n_i)/C_Q$ is the local expected gap from the Bateman--Horn density. The mean of $\hat{g}$ is $1.023$ (close to the theoretical $1$), with standard deviation $0.974$.

Under the simplest heuristic (Cram\'er model), consecutive primes of a polynomial behave like independent events with local probability $1/E_i$, predicting that $\hat{g}$ follows a standard exponential distribution with density $e^{-x}$. Table~\ref{tab:normgap} compares the empirical distribution with this prediction and with a \emph{mod-$11$ corrected} model that accounts for the discrete constraint imposed by the forbidden residues $\{4, 5, 7, 8\} \pmod{11}$.

\begin{table}[ht]
\centering
\caption{Empirical distribution of normalized gaps $\hat{g}_i = g_i / E_i$, where $E_i = (4\log n_i)/C_Q$ is the local expected gap}
\label{tab:normgap}
\smallskip
\begin{tabular}{@{}rrrrr@{}}
\toprule
Bin & Observed & Fraction & Exponential & Mod-$11$ corrected \\
\midrule
$[0,\, 0.25)$    & $988{,}981$  & $0.191$ & $0.221$ & $0.192$ \\
$[0.25,\, 0.50)$  & $835{,}415$  & $0.161$ & $0.172$ & $0.158$ \\
$[0.50,\, 0.75)$  & $860{,}533$  & $0.166$ & $0.134$ & $0.170$ \\
$[0.75,\, 1.00)$  & $485{,}622$  & $0.094$ & $0.104$ & $0.095$ \\
$[1.00,\, 1.50)$  & $821{,}858$  & $0.159$ & $0.145$ & $0.143$ \\
$[1.50,\, 2.00)$  & $481{,}813$  & $0.093$ & $0.088$ & $0.100$ \\
$[2.00,\, 3.00)$  & $457{,}348$  & $0.088$ & $0.086$ & $0.089$ \\
$[3.00,\, 4.00)$  & $158{,}962$  & $0.031$ & $0.031$ & $0.034$ \\
$[4.00,\, 6.00)$  & $77{,}769$   & $0.015$ & $0.016$ & $0.016$ \\
$\ge 6.00$       & $10{,}140$    & $0.002$ & $0.002$ & $0.003$ \\
\bottomrule
\end{tabular}
\end{table}

The corrected model is defined as follows. For a raw gap $g$ starting at a safe residue $r \bmod 11$, the transition weight $w(g) = \frac{1}{7}\#\{r \in S : (r+g) \bmod 11 \in S\}$ captures the fraction of safe-to-safe transitions, where $S = \{0,1,2,3,6,9,10\}$. Gaps $g \equiv 0 \pmod{11}$ receive weight $w = 1$, while $g \equiv 5, 6 \pmod{11}$ receive the minimum $w = 3/7$. The corrected bin probabilities are $P_{\mathrm{corr}}([a,b)) = Z^{-1}\sum_{g=\lceil aE\rceil}^{\lfloor bE\rfloor} w(g)\,e^{-g/E}/E$ with $E = \bar{g} = 19.31$ and $Z$ a normalization constant.

With $5.18 \times 10^6$ data points, classical $\chi^2$ tests are uninformative: even tiny model deviations produce enormous test statistics that always reject.  We therefore quantify the fit using sample-size-independent effect-size metrics.  The \emph{total variation distance} (TVD) and \emph{Kullback--Leibler divergence} from the observed distribution to each model are:
\begin{equation}\label{eq:effectsize}
\begin{aligned}
\mathrm{TVD}_{\mathrm{exp}} &= 0.054, &\quad D_{\mathrm{KL}}^{\mathrm{exp}} &= 7.6 \times 10^{-3}, \\
\mathrm{TVD}_{\mathrm{mod\text{-}11}} &= 0.019, &\quad D_{\mathrm{KL}}^{\mathrm{mod\text{-}11}} &= 1.6 \times 10^{-3}.
\end{aligned}
\end{equation}
The mod-$11$ correction reduces the TVD by a factor of~$2.8$ and the KL divergence by a factor of~$4.8$, confirming that $p = 11$ is the dominant source of deviation from exponential behavior.  The residual TVD of $0.019$ is attributable to the higher sieving primes $p = 31, 41, 61, \ldots$, each of which imposes additional transition-weight modulations on the gap distribution; a full multi-prime model incorporating all primes $p \equiv 1 \pmod{5}$ up to~$P$ would, in the limit $P \to \infty$, converge to the exact gap law.


% ====================================================================
\section{Consecutive-prime tuplets}
\label{sec:tuplets}
% ====================================================================

\subsection{Definitions and counts}

\begin{definition}
A \emph{$k$-tuplet} for $Q$ is a maximal block of $k$ consecutive integers $\{n, n{+}1, \ldots, n{+}k{-}1\}$ such that $Q(n), Q(n{+}1), \ldots, Q(n{+}k{-}1)$ are all prime. When counting non-maximally (allowing overlaps from longer runs), we write $T_k(N)$ for the number of starting positions $n \le N$.
\end{definition}

Table~\ref{tab:tuplets} summarizes our findings. The data for $n \le 10^8$ come from the exhaustive search described above; the data for $n \le 1.9 \times 10^9$ come from an extended sieve-based search with PFGW verification of primality for candidates near known quadruplets.

\begin{table}[ht]
\centering
\caption{$k$-tuplet counts (maximal runs)}
\label{tab:tuplets}
\smallskip
\begin{tabular}{@{}rrrr@{}}
\toprule
$k$ & $n \le 10^8$ & $n \le 1.9 \times 10^9$ & Smallest $n$ \\
\midrule
$2$ (pairs) & $254{,}833$ & --- & $2$ \\
$3$ (triplets)  & $15{,}307$  & --- & $427$ \\
$4$ (quadruplets) & $924$ & $11{,}063$ & $383$ \\
$5$ (quintuplets) & $54$ & $458$ & $1{,}915{,}329$ \\
$6$ (sextuplets)  & $3$  & $25$  & $10{,}291{,}136$ \\
$\ge 7$ & $0$ & $0$ & (impossible) \\
\bottomrule
\end{tabular}
\end{table}


\subsection{Hardy--Littlewood heuristic for $k$-tuplets}

To predict $T_k(N)$, we generalize the Bateman--Horn conjecture to the $k$-tuple of polynomials $f_j(n) = Q(n + j)$ for $j = 0, 1, \ldots, k-1$. Following the framework of Hardy and Littlewood~\cite{hardy1923}, the singular series becomes
\begin{equation}\label{eq:Sk}
\mathfrak{S}_k \;=\; \prod_{p} \frac{1 - \omega_k(p)/p}{(1 - 1/p)^k},
\end{equation}
where $\omega_k(p) = \#\{n \bmod p : Q(n+j) \equiv 0 \pmod{p} \text{ for some } 0 \le j < k\}$ is the number of ``unsafe'' starting positions modulo~$p$.

For primes $p \not\equiv 1 \pmod{5}$ (and $p \ne 5$), we have $\omega_k(p) = 0$, so the factor is $(1 - 1/p)^{-k}$. For $p \equiv 1 \pmod{5}$, the four roots of $Q \pmod{p}$ contribute: $\omega_k(p) = p - s_k(p)$, where $s_k(p)$ is the number of starting positions modulo $p$ for which the window $\{n, \ldots, n+k-1\}$ avoids all four roots.

The predicted count is
\begin{equation}\label{eq:Tkpred}
T_k(N) \;\sim\; \mathfrak{S}_k \int_2^N \frac{dt}{(\log Q(t))^k}\,.
\end{equation}

\begin{remark}\label{rem:logproduct}
More precisely, the Bateman--Horn conjecture applied to the $k$-tuple $(Q(n), Q(n{+}1), \ldots, Q(n{+}k{-}1))$ gives
\[
T_k(N) \;\sim\; \mathfrak{S}_k \int_2^N \frac{dt}{\log Q(t) \cdot \log Q(t{+}1) \cdots \log Q(t{+}k{-}1)}\,.
\]
Since $Q(t) = 5t^4(1 + O(1/t))$, we have $\log Q(t+j) = \log Q(t) + O(k/t)$ for fixed $k$, and the two integral forms agree asymptotically. At $N = 10^8$, the relative difference between $\prod_{j=0}^{k-1}\log Q(t+j)$ and $(\log Q(t))^k$ is of order $k^2/(t\log t) < 10^{-5}$ for all $t \ge 100$, so the simpler form \eqref{eq:Tkpred} introduces negligible error for the data in Table~\ref{tab:Sk}. For smaller $N$ or larger $k$, the product form would be preferable.
\end{remark}

Table~\ref{tab:Sk} compares these predictions with the observed \emph{overlapping} counts $T_k(10^8)$. Note that these differ from the maximal-run counts in Table~\ref{tab:tuplets}: for example, a single maximal $6$-tuplet at $n$ contributes one entry to the ``$k=6$'' row of Table~\ref{tab:tuplets}, but contributes to $T_6$, $T_5$, $T_4$, $T_3$, and $T_2$ in Table~\ref{tab:Sk} (providing respectively $1$, $2$, $3$, $4$, and $5$ overlapping sub-tuplets).

\begin{table}[ht]
\centering
\caption{Singular series and tuplet predictions at $N = 10^8$ (overlapping counts $T_k$)}
\label{tab:Sk}
\smallskip
\begin{tabular}{@{}rrrrc@{}}
\toprule
$k$ & $\mathfrak{S}_k$ & $\displaystyle\int_2^{10^8}\!\frac{dt}{(\log Q)^k}$ & Predicted & Observed \\
\midrule
$2$ & $14.42$ & $19{,}903$ & $286{,}951$ & $288{,}450$ \\
$3$ & $60.81$ & $282.9$ & $17{,}200$  & $17{,}329$ \\
$4$ & $255.0$ & $4.050$ & $1{,}033$  & $1{,}041$ \\
$5$ & $903.2$ & $0.0594$ & $53.6$     & $60$ \\
$6$ & $2{,}220$ & $0.0011$ & $2.4$     & $3$ \\
\bottomrule
\end{tabular}
\end{table}

For $k = 2, 3, 4$, the predictions agree with the data to within $1\%$, providing strong evidence for the Hardy--Littlewood model. The $k = 5, 6$ cases show larger deviations ($12\%$ and $26\%$ respectively), as expected from the small sample sizes.

\begin{remark}
The vanishing of $\mathfrak{S}_7$ provides an independent, heuristic-free confirmation of Theorem~\ref{thm:7tuplet}: at $p = 11$, we have $s_7(11) = 0$, so $\omega_7(11) = 11$, and the factor $(1 - 11/11)/(1 - 1/11)^7 = 0$ forces $\mathfrak{S}_7 = 0$.
\end{remark}


\subsection{Distribution of sextuplets}

The $25$ sextuplets found up to $n \approx 1.9 \times 10^9$ begin at the following values of $n$:
\begin{gather*}
10{,}291{,}136,\;\; 45{,}197{,}062,\;\; 89{,}166{,}471,\;\; 139{,}663{,}665,\;\; 144{,}073{,}631,\;\; 150{,}737{,}728, \\
159{,}933{,}002,\;\; 169{,}996{,}638,\;\; 181{,}637{,}322,\;\; 226{,}237{,}427,\;\; 233{,}892{,}349,\;\; 308{,}374{,}240, \\
506{,}989{,}173,\;\; 548{,}169{,}675,\;\; 570{,}232{,}177,\;\; 608{,}369{,}804,\;\; 843{,}734{,}362,\;\; 865{,}999{,}594, \\
1{,}055{,}031{,}536,\;\; 1{,}123{,}085{,}445,\;\; 1{,}188{,}401{,}113,\;\; 1{,}393{,}470{,}692,\;\; 1{,}625{,}095{,}965, \\
1{,}686{,}424{,}793,\;\; 1{,}932{,}852{,}337.
\end{gather*}
A sextuplet at $n$ requires $n, n{+}1, \ldots, n{+}5$ to all be safe modulo~$11$. Since the unique maximal block of $6$ consecutive safe residues is $\{9, 10, 0, 1, 2, 3\}$, every sextuplet must satisfy $n \equiv 9 \pmod{11}$. We have verified that all $25$ values listed above do satisfy this necessary condition.


% ====================================================================
\section{Large probable prime generation}
\label{sec:large}
% ====================================================================

\subsection{Three-stage pipeline}

The arithmetic rigidity of $Q(n)$ (Corollary~\ref{cor:divisibility}) makes it an excellent candidate for generating titanic probable primes ($\ge 1{,}000$ digits) and larger. We developed a three-stage pipeline:

\medskip
\noindent\textbf{Stage~1: Modular sieve (C++).}
For a chosen base $n_0$ (e.g., $n_0 = 10^{2499}$ to target $10{,}000$-digit values of $Q$), we sieve the interval $\{n_0, n_0{+}1, \ldots, n_0{+}R-1\}$ for offsets $i$ such that $Q(n_0 + i)$ has no prime factor $p \le B$ with $p \equiv 1 \pmod{5}$.

For each such prime $p$, we compute the four roots $r_1, \ldots, r_4$ of $Q \pmod{p}$ using a primitive fifth root of unity, then mark all $i$ with $n_0 + i \equiv r_j \pmod{p}$. With $B = 10^9$ and $R = 1{,}200{,}000$, this eliminates approximately $99.97\%$ of candidates, leaving several thousand ``survivors.''

The sieve is implemented in C++ using $128$-bit arithmetic for modular reductions (cf.~\cite{pritchard1987,sorenson1998} for background on prime-number sieves) and runs in approximately one hour on a single core (AMD Ryzen 9 5950X, 3.4\,GHz).

\medskip
\noindent\textbf{Stage~2: Target construction (Python).}
The surviving offsets are expanded into full decimal representations of $Q(n_0 + i)$, producing numbers of $\sim 4d$ digits where $d$ is the number of digits in $n_0$. These are formatted for input to PFGW.

\medskip
\noindent\textbf{Stage~3: PRP testing (PFGW).}
Each target is subjected to a strong Fermat probable-prime test using PFGW~\cite{pfgw}, with base~$3$ and a subsequent Brillhart--Lehmer--Selfridge (BLS) partial factorization test when applicable. A $10{,}000$-digit candidate requires approximately $30$ seconds; a $100{,}000$-digit candidate requires several hours; a $200{,}000$-digit candidate requires roughly one day.

\begin{remark}
The numbers produced by the pipeline above are initially \emph{probable primes} (PRPs) based on a strong Fermat test.  To obtain a rigorous primality proof, we subjected the smallest candidate $Q(10^{2499} + 3{,}469)$ to the ECPP algorithm~\cite{atkinmorain1993} using the implementation Primo~4.3.3~\cite{primo}. The computation required $286{,}269$ seconds of wall-clock time ($\approx 3.3$ days; $43.4$ CPU-days across multiple cores), producing a certificate chain of $887$ steps that unconditionally proves $Q(10^{2499} + 3{,}469)$ is prime. This $9{,}997$-digit number is the largest proven prime of the form $n^5 - (n-1)^5$ currently listed in the Prime Pages database~\cite{primepages}. The Primo certificate file is included in the supplementary materials. The remaining $174$ candidates in Appendix~\ref{app:offsets}, as well as the $50{,}000$-, $100{,}000$-, and $200{,}000$-digit PRPs, await certification.
\end{remark}

\subsection{Results}

Table~\ref{tab:largePRP} summarizes the probable primes found to date.

\begin{table}[ht]
\centering
\caption{Probable primes (PRPs) of the form $Q(n) = n^5 - (n-1)^5$}
\label{tab:largePRP}
\smallskip
\begin{tabular}{@{}rllrl@{}}
\toprule
Digit class & Base $n_0$ & Offset range & PRPs & Status \\
\midrule
$\sim 10{,}000$ & $10^{2499}$ & $3{,}469 \le k \le 1{,}162{,}743$ & $175$ & $k=3{,}469$ certified prime \\
$\sim 50{,}000$ & $10^{12499}$ & $k \in \{5{,}730,\; 29{,}658\}$ & $2$ & PRP \\
$\sim 100{,}000$ & $10^{24999}$ & $k = 25{,}098$ & $1$ & PRP \\
$\sim 200{,}000$ & $10^{49999}$ & $k = 7{,}506{,}520$ & $1$ & PRP \\
\bottomrule
\end{tabular}
\end{table}

In the $10{,}000$-digit search, the sieve covered offsets $0 \le i < 1{,}200{,}000$ from the base $n_0 = 10^{2499}$, with sieve depth $B = 10^9$. Subsequent PFGW testing of all survivors identified $175$ probable primes $Q(10^{2499}+k)$, each having $9{,}997$ decimal digits. The complete list of all $175$ offsets is provided in Appendix~\ref{app:offsets}. As described in Remark~7, the smallest of these---the $9{,}997$-digit number
\[
Q(10^{2499} + 3{,}469) \;=\; (10^{2499} + 3{,}469)^5 - (10^{2499} + 3{,}468)^5
\]
---has been rigorously certified prime via ECPP.

The mean spacing between consecutive PRPs is $\approx 6{,}662$, consistent with the Bateman--Horn density $C_Q / \!\log Q(n_0) \approx 1.6 \times 10^{-4}$. We note the close pair $k = 106{,}459$ and $k = 106{,}520$ (gap $61$), suggesting that tuplet-like clustering persists even at the $10{,}000$-digit scale.

For the $50{,}000$-digit class, two probable primes emerged from a similar pipeline. The $100{,}000$-digit probable prime $Q(10^{24999} + 25{,}098)$ has $99{,}997$ decimal digits. Most recently, a search at base $n_0 = 10^{49999}$ yielded the $200{,}000$-digit probable prime
\[
Q(10^{49999} + 7{,}506{,}520) \;=\; (10^{49999} + 7{,}506{,}520)^5 - (10^{49999} + 7{,}506{,}519)^5
\]
with $199{,}997$ decimal digits, which is the largest known (probable) prime of the form $n^5 - (n-1)^5$ as of March 2026~\cite{primepages}.


% ====================================================================
\section{Concluding remarks}
\label{sec:conclusion}
% ====================================================================

The polynomial $Q(n) = n^5 - (n-1)^5$ displays a remarkably clean arithmetic structure rooted in the fifth cyclotomic field. The restriction that all prime divisors satisfy $p \equiv 1 \pmod{5}$ boosts the prime density by a factor of $C_Q \approx 3.678$, enables the existence of consecutive-prime sextuplets, and simultaneously imposes the sharp bound $k \le 6$ on tuplet length. All quantitative predictions from the Bateman--Horn and Hardy--Littlewood frameworks are confirmed by the data to high precision.

Several natural directions for further work include:
\begin{enumerate}[label=(\roman*)]
\item \emph{Further primality certifications.} Having certified $Q(10^{2499} + 3{,}469)$ as a $9{,}997$-digit proven prime (Remark~7), we plan to certify additional candidates from Appendix~\ref{app:offsets}. Each ECPP run at this scale requires $\approx 3$ days of wall-clock time (Primo~4.3.3 on a multi-core workstation), making batch certification of all $175$ candidates a significant but tractable computational project.

\item \emph{Megaprime search.} The sieve pipeline scales well: replacing the base $n_0 = 10^{2499}$ with $n_0 = 10^{249999}$ would target $10^6$-digit probable primes. The sieve cost is dominated by the number of active primes $p \le B$ with $p \equiv 1 \pmod{5}$, which is independent of $n_0$; only the PRP testing cost grows (roughly as $d^2$ for $d$-digit numbers under FFT-based arithmetic~\cite{crandall2005}).

\item \emph{Higher cyclotomic analogues.} For odd $k$, the polynomial $n^k - (n-1)^k$ is governed by $\Phi_k$, and the sieving restriction is $p \equiv 1 \pmod{k}$. The cases $k = 7, 11, 13$ would give even stronger density boosts ($C_f$ grows with $k$) but require deeper sieving and present more complex root structures.

\item \emph{Tuplet asymptotics.} The singular series $\mathfrak{S}_k$ for $Q$ grow rapidly with $k$ (Table~\ref{tab:Sk}), but the integral factor $\int dt/(\log Q)^k$ shrinks faster. A careful analysis of when $\mathfrak{S}_k \int dt/(\log Q)^k$ drops below~$1$ (expected number $< 1$) would give an effective ``last expected $k$-tuplet'' bound.

\item \emph{CRT-based targeted search.} Rather than sieving a contiguous interval, one could use the Chinese Remainder Theorem to solve a system of congruences $n \not\equiv r_j \pmod{p}$ simultaneously for many primes $p$, directly computing candidate $n$-values that survive deep sieving. This approach, if combined with lattice reduction, could accelerate the search for tuplets or large primes by orders of magnitude. We leave this as an open problem.
\end{enumerate}


% ====================================================================
\appendix
\section{Complete list of 10,000-digit PRP offsets}
\label{app:offsets}
% ====================================================================

The following is the complete list of all $175$ offsets $k$ ($3{,}469 \le k \le 1{,}162{,}743$) for which $Q(10^{2499}+k)$ is a probable prime, as described in Section~\ref{sec:large}:
\begin{gather*}
3469,\; 6073,\; 14791,\; 15367,\; 20398,\; 26943,\; 44759,\; 48337,\; 49436,\; 50368,\; 60716, \\
61899,\; 64870,\; 65769,\; 71174,\; 73025,\; 87603,\; 88183,\; 91080,\; 93129,\; 94912, \\
100905,\; 105403,\; 106459,\; 106520,\; 109353,\; 115504,\; 117292,\; 131871,\; 135981, \\
137534,\; 138028,\; 139571,\; 142308,\; 155852,\; 156578,\; 162669,\; 179711,\; 182412, \\
185742,\; 187007,\; 190434,\; 203823,\; 205594,\; 207126,\; 208795,\; 214215,\; 216989, \\
223717,\; 229791,\; 230914,\; 235979,\; 238029,\; 242233,\; 242396,\; 243111,\; 245819, \\
261320,\; 262346,\; 263803,\; 266465,\; 281890,\; 283375,\; 289677,\; 293729,\; 300741, \\
301347,\; 307443,\; 311049,\; 325149,\; 328230,\; 328680,\; 333465,\; 334550,\; 335280, \\
351457,\; 364478,\; 384056,\; 384285,\; 387584,\; 392416,\; 393658,\; 411127,\; 418036, \\
419267,\; 426224,\; 429397,\; 437129,\; 444198,\; 445412,\; 445939,\; 483770,\; 495067, \\
496221,\; 532085,\; 541966,\; 558371,\; 559053,\; 568659,\; 575761,\; 578951,\; 587286, \\
595302,\; 598946,\; 604496,\; 615156,\; 620279,\; 621807,\; 623113,\; 644365,\; 646958, \\
650606,\; 650928,\; 676413,\; 677717,\; 681672,\; 684941,\; 696148,\; 709895,\; 722259, \\
745011,\; 747585,\; 747809,\; 758768,\; 788690,\; 789870,\; 796358,\; 796935,\; 802798, \\
805237,\; 812669,\; 817212,\; 817904,\; 819548,\; 822249,\; 830291,\; 838259,\; 846837, \\
849237,\; 851349,\; 860556,\; 865248,\; 874512,\; 891826,\; 903335,\; 919083,\; 922789, \\
923219,\; 924617,\; 931960,\; 941218,\; 946285,\; 968134,\; 975820,\; 984051,\; 988826, \\
994495,\; 996189,\; 1000966,\; 1008789,\; 1017419,\; 1026127,\; 1029167,\; 1055772, \\
1057771,\; 1072702,\; 1084128,\; 1089763,\; 1113852,\; 1128875,\; 1134903,\; 1152880, \\
1153151,\; 1161083,\; 1162743.
\end{gather*}


% ====================================================================
% ACKNOWLEDGMENTS
% ====================================================================

\medskip
\noindent\textbf{Acknowledgments.}
The computations were performed on infrastructure at GUT Geoservice Inc., Montr\'eal. The author thanks the developers of PFGW for making their software freely available.

\medskip
\noindent\textbf{Code availability.}
The C++ sieve, Python pipeline scripts, and all raw data files used in this paper are publicly available at \url{https://github.com/GUTgeoservice/Q5-primes}.

% ====================================================================
% REFERENCES
% ====================================================================

\begin{thebibliography}{99}

\bibitem{atkinmorain1993}
A.\,O.\,L.~Atkin and F.~Morain,
\emph{Elliptic curves and primality proving},
Math.\ Comp.\ \textbf{61} (1993), 29--68.

\bibitem{bateman1962}
P.\,T.~Bateman and R.\,A.~Horn,
\emph{A heuristic asymptotic formula concerning the distribution of prime numbers},
Math.\ Comp.\ \textbf{16} (1962), 363--367.

\bibitem{cohen2007}
H.~Cohen,
\emph{Number Theory, Volume~II: Analytic and Modern Tools},
Graduate Texts in Math., vol.~240, Springer, New York, 2007.

\bibitem{crandall2005}
R.~Crandall and C.~Pomerance,
\emph{Prime Numbers: A Computational Perspective},
2nd ed., Springer, New York, 2005.

\bibitem{dubner2002}
H.~Dubner, T.~Forbes, N.~Lygeros, M.~Mizony, H.~Nelson, and P.~Zimmermann,
\emph{Ten consecutive primes in arithmetic progression},
Math.\ Comp.\ \textbf{71} (2002), 1323--1330.

\bibitem{galway2004}
W.~Galway,
\emph{Analytic computation of the prime-counting function},
Ph.\,D.\ thesis, University of Illinois at Urbana-Champaign, 2004.

\bibitem{hardy1923}
G.\,H.~Hardy and J.\,E.~Littlewood,
\emph{Some problems of `Partitio Numerorum'; III: On the expression of a number as a sum of primes},
Acta Math.\ \textbf{44} (1923), 1--70.

\bibitem{iwaniec2004}
H.~Iwaniec and E.~Kowalski,
\emph{Analytic Number Theory},
Amer.\ Math.\ Soc.\ Colloq.\ Publ., vol.~53, Amer.\ Math.\ Soc., Providence, RI, 2004.

\bibitem{languasco2010}
A.~Languasco and A.~Zaccagnini,
\emph{Computing the Mertens and Meissel--Mertens constants for sums over arithmetic progressions},
Experiment.\ Math.\ \textbf{19} (2010), 279--284.

\bibitem{oliveira2014}
T.~Oliveira~e~Silva, S.~Herzog, and S.~Pardi,
\emph{Empirical verification of the even Goldbach conjecture and computation of prime gaps up to $4 \cdot 10^{18}$},
Math.\ Comp.\ \textbf{83} (2014), 2033--2060.

\bibitem{pfgw}
OpenPFGW Project,
\emph{PFGW --- a primality testing program},
\url{https://sourceforge.net/projects/openpfgw/}.

\bibitem{primepages}
C.~Caldwell,
\emph{The Prime Pages: prime number research, records and resources},
\url{https://t5k.org/}, accessed March 2026.

\bibitem{primo}
M.~Martin,
\emph{Primo --- a primality proving program based on ECPP},
\url{https://www.ellipsa.eu/}, 2023.

\bibitem{pritchard1987}
P.~Pritchard,
\emph{Linear prime-number sieves: A family tree},
Sci.\ Comput.\ Programming \textbf{9} (1987), 17--35.

\bibitem{sorenson1998}
J.~Sorenson,
\emph{Trading time for space in prime number sieves},
in: \emph{Algorithmic Number Theory} (ANTS-III), Lecture Notes in Comput.\ Sci., vol.~1423, Springer, 1998, pp.~179--195.

\bibitem{sorenson2024}
J.~Sorenson,
\emph{The Bateman--Horn conjecture: heuristics, history, and applications},
in: \emph{Computational and Analytic Number Theory} (A.\,Granville and Z.\,Rudnick, eds.), Fields Inst.\ Commun., to appear.

\end{thebibliography}

\end{document}